\documentclass[11pt]{article}
\usepackage[letterpaper,margin=0.82in]{geometry}
\usepackage{amsmath,amssymb,amsthm,mathtools}
\usepackage{graphicx}
\usepackage{microtype}
\usepackage[hidelinks]{hyperref}

\newtheorem{theorem}{Theorem}
\newtheorem{lemma}{Lemma}
\newcommand{\Q}{\mathbb{Q}}

\title{All-order parity of the Euler-correction\\regular C-fraction}
\author{Automated mathematics run \texttt{fa\_banach\_001}\\Lane 3 --- candidate for human review}
\date{17 August 2026}

\begin{document}
\maketitle

\noindent\textbf{Status.}
Candidate full solution, likely valid.  The result gives an affirmative answer
to the open problem after Theorem~1 of Hongmin Xu and Xu You,
\emph{Inequalities for Euler--Mascheroni constant}, arXiv:1407.3865
\cite{XuYou}.

\section{Source problem}

For coefficients \(a_j\), the source uses
\[
 R_k(n)=\cfrac{a_1}{
 n+\cfrac{a_2n}{
 n+\cfrac{a_3n}{
 n+\ddots+\cfrac{a_{k-1}n}{n+a_k}}}},                   \tag{1}
\]
and \(r_k(n)=H_n-\log n-R_k(n)\),
choosing the coefficients successively to make \(r_k(n)-\gamma\) converge as
fast as possible.  Its computations through \(a_{13}\) suggest the following
statement, printed on source-PDF page 3.

\begin{center}
 \includegraphics[width=0.96\linewidth]{figures/open_problem_crop.png}
\end{center}

\noindent\textbf{Open problem (transcribed).}
For every \(k\geq1\), is \(a_{2k+1}=-a_{2k}\)?

\section{Proof intuition}

After putting \(x=1/n\), (1) becomes an ordinary regular C-fraction.
The Euler correction has an asymptotic series consisting of \(x/2\) plus
\emph{only even powers}.  The nonlinear C-fraction tail algorithm converts
this parity into pairs: extracting one coefficient from a tail of the form
\(\alpha x+E(x^2)\) leaves a new tail whose next linear coefficient is the
negative of the one just extracted.

Parity alone would not rule out the algorithm terminating.  The even
coefficients are absolute Bernoulli numbers, however, and form a strict
Stieltjes moment sequence with an explicit positive density.  Consequently
all relevant Hankel determinants are positive and the Stieltjes fraction,
hence the paired C-fraction, continues nondegenerately to every order.

\section{Formal C-fraction preliminaries}

For a formal series \(A(x)=O(x)\), define its regular C-fraction tails by
\[
 T_1=A,
 \qquad a_j=[x]T_j,
 \qquad T_{j+1}=\frac{a_jx}{T_j}-1.                       \tag{2}
\]
Whenever \(a_j\neq0\), this is equivalent to
\[
 T_j=\frac{a_jx}{1+T_{j+1}}.
\]
Iteration gives
\[
 A(x)=\cfrac{a_1x}{1+\cfrac{a_2x}{1+\cfrac{a_3x}{1+\ddots}}}.          \tag{3}
\]
Truncating after \(a_k\) changes the series first at order \(x^{k+1}\).
Thus (2) uniquely produces the coefficient at each stage that gives the
fastest \(k\)-level approximation, provided the next linear coefficient is
nonzero.

\begin{lemma}[parity mechanism]\label{lem:parity}
Suppose an odd-indexed tail has the form
\[
 T(x)=\alpha x+x^2C(x^2),\qquad \alpha C(0)\neq0.
\]
The next two coefficients extracted by (2) are opposites.  The tail after
those two extractions again has the form linear term plus an even series.
\end{lemma}

\begin{proof}
The first new tail is
\[
 U(x)=\frac{\alpha x}{T(x)}-1
     =-\frac{xC(x^2)}{\alpha+xC(x^2)}.
\]
Hence its linear coefficient is \(b=-C(0)/\alpha\).  Extracting it gives
\[
 V(x)=\frac{bx}{U(x)}-1
     =-bx+\left(\frac{C(0)}{C(x^2)}-1\right).              \tag{4}
\]
The second coefficient is therefore \(-b\), and the parenthesized term is an
even series with zero constant term.
\end{proof}

\section{The all-order result}

Let \(B_j\) denote the Bernoulli numbers in the convention \(B_2=1/6\).
Euler--Maclaurin gives the formal asymptotic expansion
\[
 H_n-\log n-\gamma
 \sim \frac{1}{2n}-\sum_{m\geq1}\frac{B_{2m}}{2m\,n^{2m}}.
\]
With \(x=1/n\), write this as
\[
 F(x)=\frac{x}{2}-x^2S(x^2),
 \qquad
 S(z)=\sum_{m\geq0}(-1)^m\mu_m z^m,
 \qquad
 \mu_m=\frac{|B_{2m+2}|}{2m+2}.                           \tag{5}
\]

\begin{theorem}\label{thm:main}
The fastest-approximation coefficients in (1) exist uniquely to every order.
They satisfy
\[
 a_{2k}>0,
 \qquad
 \boxed{a_{2k+1}=-a_{2k}}
 \qquad(k\geq1).
\]
In particular, the source open problem has an affirmative answer.
\end{theorem}

\begin{proof}
The classical formula for even Bernoulli numbers gives
\begin{align*}
 \mu_m
 &=\frac{2(2m+1)!\,\zeta(2m+2)}{(2\pi)^{2m+2}}\\
 &=\int_0^\infty
    \frac{t^m}{e^{2\pi\sqrt t}-1}\,dt.                    \tag{6}
\end{align*}
For the second equality, substitute \(t=u^2\) and use the standard Gamma--zeta
integral.  The density in (6) is strictly positive on \((0,\infty)\).
Consequently, for every nonzero polynomial \(P\),
\[
 \int_0^\infty |P(t)|^2\frac{dt}{e^{2\pi\sqrt t}-1}>0,
 \qquad
 \int_0^\infty t|P(t)|^2\frac{dt}{e^{2\pi\sqrt t}-1}>0.  \tag{7}
\]
Thus every ordinary and shifted Hankel moment matrix of \((\mu_m)\) is
positive definite.

The formal Stieltjes continued-fraction theorem now applies: strict positivity
in (7) yields
\[
 S_0(z):=S(z)=\frac{c_0}{1+zS_1(z)},
 \qquad
 S_j(z)=\frac{c_j}{1+zS_{j+1}(z)}\quad(j\geq1),            \tag{8}
\]
with every \(c_j=S_j(0)>0\)~\cite[Ch.~3]{Wall}.  One may also see (8)
directly from Gram--Schmidt: the denominators are the orthogonal polynomials
of the positive measure in (6), and their squared norms and recurrence
coefficients cannot vanish.

We now run (2) on \(F\).  Put \(\alpha=1/2\) and
\(b_1=c_0/\alpha=1/6\).  Directly from (5) and (8),
\begin{align*}
 T_2&=\frac{xS_0(x^2)}{\alpha-xS_0(x^2)},\\
 T_3&=-b_1x+x^2S_1(x^2).
\end{align*}
Thus \(a_2=b_1\) and \(a_3=-b_1\).  Inductively, if
\[
 T_{2j+1}=-b_jx+x^2S_j(x^2),
\]
define \(b_{j+1}=c_j/b_j>0\).  Using (8) in the same two algebraic
extractions as in Lemma~\ref{lem:parity} gives
\[
 a_{2j+2}=b_{j+1},
 \qquad a_{2j+3}=-b_{j+1},
 \qquad
 T_{2j+3}=-b_{j+1}x+x^2S_{j+1}(x^2).                      \tag{9}
\]
All extractions are nondegenerate because \(b_j,c_j>0\).  Hence (9) holds for
all \(j\), proving both existence and the asserted signs.

Finally, normalizing (1) by \(x=1/n\) turns its \(k\)-level fraction exactly
into the \(k\)-th convergent of (3).  The nonzero next coefficient makes its
error have exact order \(x^{k+1}\); any different choice at the first differing
stage loses that order.  Therefore these are precisely the source's uniquely
fastest coefficients.
\end{proof}

\section{Constructive recurrence and exact verification}

Equations (8)--(9) give a constructive description.  If \(c_j\) are the
positive Stieltjes-fraction constants of the moment sequence (6), then
\[
 b_1=\frac16,
 \qquad b_{j+1}=\frac{c_j}{b_j},
 \qquad a_{2j}=b_j,
 \qquad a_{2j+1}=-b_j.                                    \tag{10}
\]
Because all moments are rational, every \(a_j\) is rational.  For example,
continuing beyond the source computation gives
\[
 a_{14}=\frac{3754087889491759}{2440840521848406},
 \qquad a_{15}=-a_{14}.
\]

The exact SymPy script \texttt{code/check\_cfraction\_parity.py} implements
(2) on (5).  It reproduces every coefficient printed by the source through
\(a_{13}\), verifies the twelve complete pairs available through \(a_{25}\)
while computing through \(a_{26}\), and confirms the displayed new values.
This finite calculation is QA only; Theorem~\ref{thm:main} is the all-order
proof.

\section{Novelty check, limitations, and review focus}

A bounded search on 17 August 2026 checked the run indexes; exact web and
arXiv queries for the displayed coefficient relation and the source title;
the source's indexed citing-work list; and later continued-fraction papers by
the same circle of authors.  In particular, Cao's later paper cites the source
and develops a different multiple-correction fraction, but still reports only
the coefficients through \(a_{13}\) from (1)~\cite{Cao}.  No paper proving the
all-order opposite-pair law was found.  This supports, but does not establish,
novelty; the search was not an exhaustive MathSciNet or zbMATH review.

The theorem settles the source's coefficient conjecture and proves the
coefficients exist at every formal asymptotic order.  It does not assert
convergence of the resulting infinite C-fraction for fixed \(n\), which the
source question does not require.  Human review should focus on the standard
Stieltjes-fraction implication in (8) and on the equivalence between the
source's ``fastest'' recursion and regular C-fraction coefficient matching.

\begin{thebibliography}{9}
\bibitem{XuYou}
Hongmin Xu and Xu You,
\emph{Inequalities for Euler--Mascheroni constant},
arXiv:1407.3865 (2014); published as
\emph{Continued fraction inequalities for Euler--Mascheroni constant},
J. Inequal. Appl. 2014:343.

\bibitem{Wall}
H.~S. Wall,
\emph{Analytic Theory of Continued Fractions},
D. Van Nostrand, 1948, Chapter 3.

\bibitem{Cao}
Xiaodong Cao,
\emph{Multiple-Correction and Continued Fraction Approximation},
arXiv:1410.2953; J. Math. Anal. Appl. 424 (2015), 1425--1446.
\end{thebibliography}

\end{document}
